\documentclass[a4paper,11pt]{article}
\usepackage[utf8]{inputenc}
\usepackage[T1]{fontenc}
\usepackage[a4paper,margin=2.5cm]{geometry}
\usepackage{amsmath,amssymb,amsfonts}
\usepackage{graphicx}
\usepackage{booktabs}
\usepackage{hyperref}
\usepackage{natbib}
\usepackage{xcolor}
\usepackage{microtype}
\usepackage{siunitx}

\hypersetup{
  colorlinks=true,
  linkcolor=blue!60!black,
  citecolor=blue!60!black,
  urlcolor=blue!60!black
}

% ── Macros ──────────────────────────────────────────────────────
\newcommand{\DC}{d_c}
\newcommand{\Geff}{G_{\rm eff}}
\newcommand{\rdrag}{r_{\rm drag}}
\newcommand{\sig}{\sigma_8}
\newcommand{\LCDM}{\Lambda\mathrm{CDM}}
\newcommand{\BuP}{\textsc{BuP}}
\newcommand{\camb}{\textsc{camb}}

% ── Title ────────────────────────────────────────────────────────
\title{\textbf{Contraintes cosmologiques sur une dimension spatiale\\
émergente~: dimension variable, énergie noire et $\sigma_8$\\
dans le cadre Bottom-Up Quantum Gravity}}

\author{Farid Hamdad\\[4pt]
\small Framework \BuP\ --- Bottom-Up Quantum Gravity\\
\small \url{https://github.com/Farid-Hamdad/Bottom-up-Quantum-Gravity}}

\date{Avril 2026}

\begin{document}
\maketitle

% ── Abstract ────────────────────────────────────────────────────
\begin{abstract}
Nous étudions les implications cosmologiques d'une dimension spatiale
effective~$d$ émergeant de l'intrication quantique, dans le cadre
de l'approche \textit{bottom-up} de la gravité quantique (\BuP).
À l'aide d'une version modifiée du code Boltzmann \camb, nous montrons
que toute déviation constante~$d\neq 3$ est doublement exclue~:
pour $d=2.65$, l'horizon sonore~$\rdrag$ chute de~15\% et $\sig$
augmente de~62\%, tous deux incompatibles avec les données actuelles.
Motivé par cette contrainte, nous introduisons un modèle à dimension
variable tardive fondé sur une variable intermédiaire~$X(z)$
représentant le rapport effectif entre intrication non-locale et locale.
Le profil résultant $d(z)=\DC - \Delta d\,/\,(1+X(z)^\alpha)$,
avec $\alpha=1.78$ fixé par les simulations BuP sur $N=6$--$16$
qubits, améliore significativement les ajustements aux données
DESI~2024 BAO et Pantheon+ ($\Delta\mathrm{AIC}=-6.23$,
$\Delta\mathrm{BIC}=-3.71$ par rapport à $\LCDM$).
En implémentant le couplage gravitationnel effectif
$\Geff(z)=\frac{2}{d(z)-1}\,G$ de façon cohérente dans les
équations de perturbation, nous obtenons $\sig=0.772$, en accord avec
KiDS-1000 à $0.3\sigma$. Le mécanisme est géométrique~: la légère
sur-dimensionnalité de l'Univers primordial ($\DC=3.0598>3$) inhibe
la croissance des structures à haut redshift et compense le renforcement
gravitationnel tardif. La légère sur-dimensionnalité de l'Univers primordial
($\DC=3.0598>3$) constitue le mécanisme géométrique central du cadre BuP~:
elle préserve les observables du CMB, modifie~$\Geff$ et résout la tension~$\sig$
en une seule valeur, sans paramètre libre supplémentaire.
Une cohérence multi-échelle entre simulations
($d_s\approx 2.61\pm 0.03$), échelles galactiques SPARC ($2.644\pm 0.295$)
et données cosmologiques ($2.40$--$2.71$), sur douze ordres de grandeur,
constitue le résultat observationnel le plus fort de cette analyse.
\end{abstract}

\vspace{0.5cm}
\noindent\textbf{Mots-clés~:} gravité quantique émergente, dimension effective,
intrication quantique, énergie noire, tension~$\sigma_8$, BAO, CMB, DESI~2024.

% ──────────────────────────────────────────────────────────────
\section{Introduction}
% ──────────────────────────────────────────────────────────────

La nature de l'espace-temps à un niveau fondamental reste une
question ouverte de la physique théorique. Dans les approches
\textit{bottom-up} de la gravité quantique (\BuP), la géométrie
n'est pas un fond fixe mais émerge de la structure d'intrication
d'un état quantique global~\cite{hamdad2026bup}. La dimension
spatiale effective~$d$ est alors une propriété émergente, dépendant
de l'échelle et potentiellement de l'époque cosmologique.

Des simulations numériques sur des systèmes de $N=6$--$16$ qubits
montrent que le Laplacien d'intrication conduit à une dimension
spectrale $d_s\approx 2.61\pm 0.03$, remarquablement stable en
taille de système et en couplage non-local~$\lambda$. Cette valeur
est cohérente avec la dimension effective galactique mesurée sur
l'échantillon SPARC ($d_{\rm eff}=2.644\pm 0.295$, 175~galaxies)
et avec la plage cosmologique $d(z=0)\approx 2.4$--$2.7$, sur
six ordres de grandeur en longueur.

Le modèle standard de cosmologie ($\LCDM$) fait face à des tensions
observationnelles persistantes~: la tension $H_0$ dépasse
$6\sigma$~\cite{tensions2025}, la tension $S_8/\sigma_8$ reste
$2$--$3\sigma$ selon les relevés. Ces tensions motivent l'exploration
de modifications à la physique de fond, notamment via une dimension
spatiale effective dépendant du redshift. Nous montrons que la valeur
critique $\DC=3.0598>3$, déterminée par la dynamique microscopique
du réseau d'intrication, joue un rôle unificateur~: elle garantit la
compatibilité CMB, modifie le couplage gravitationnel effectif, et
résout naturellement la tension~$\sig$ --- trois propriétés d'une seule
constante, sans ajustement cosmologique.

Cet article s'organise comme suit. La Section~\ref{sec:camb}
présente les contraintes sur $d$ constant. La Section~\ref{sec:models}
introduit le modèle~$X(z)$ et ses résultats cosmologiques.
La Section~\ref{sec:sigma8} traite la tension~$\sigma_8$ via
l'implémentation cohérente de~$\Geff$. La Section~\ref{sec:discussion}
discute les implications et les perspectives.

% ──────────────────────────────────────────────────────────────
\section{Contraintes sur $d$ constant}
\label{sec:camb}
% ──────────────────────────────────────────────────────────────

\subsection{Implémentation dans \camb}

Nous modifions \camb~1.6.5 (source Fortran) en introduisant~$d$
comme paramètre explicite dans l'équation de Friedmann à $d$
dimensions spatiales~:
\begin{equation}
  H^2(z) = \frac{8\pi G}{d(d-1)}
  \bigl[\rho_m(z) + \rho_r(z) + \rho_\Lambda^{\rm eff}(z)\bigr],
  \label{eq:friedmann}
\end{equation}
avec les lois de dilution $\rho_m\propto(1+z)^d$ et
$\rho_r\propto(1+z)^{d+1}$. Le couplage gravitationnel effectif
est $\Geff(d)=\frac{2}{d-1}\,G$ (équation de Poisson en $d$
dimensions). La validation à $d=3.0$ reproduit $\LCDM$ à
mieux de~$0.001\%$ sur~$\theta_s$ et~$\rdrag$.

\subsection{Résultats}

Le Tableau~\ref{tab:scan} et la Figure~\ref{fig:scan} présentent
les résultats du scan $d\in[2.3,\,3.1]$ à paramètres Planck~2018
fixés ($H_0=67.36$, $\Omega_bh^2=0.02237$, $\Omega_ch^2=0.1200$,
$n_s=0.965$, $A_s=2.1\times10^{-9}$, $\tau=0.054$).

\begin{table}[htbp]
\centering
\caption{Observables \camb\ en fonction de la dimension~$d$
(paramètres Planck~2018 fixés). La ligne $d=3.0$ correspond à
$\LCDM$ standard.}
\label{tab:scan}
\smallskip
\begin{tabular}{ccccccc}
\toprule
$d$ & $\theta_s$ & $\rdrag$ [Mpc] & $\sig$ & $z_{\rm eq}$ \\
\midrule
2.30 & 0.010331 & 103.2 & 2.140 & 3402.9 \\
2.50 & 0.010360 & 115.8 & 1.626 & 3402.9 \\
2.65 & 0.010378 & 125.2 & 1.312 & 3402.9 \\
2.80 & 0.010394 & 134.6 & 1.062 & 3402.9 \\
3.00 & 0.010412 & 147.1 & 0.811 & 3402.9 \\
3.05 & 0.010416 & 150.2 & 0.761 & 3402.9 \\
\bottomrule
\end{tabular}
\end{table}

\begin{figure}[htbp]
\centering
\includegraphics[width=\textwidth]{fig1_camb_scan}
\caption{Impact de la dimension effective $d$ sur trois observables
cosmologiques (paramètres Planck~2018 fixés).
\textit{Gauche}~: $\sig$ en fonction de~$d$~; la bande verte
indique la plage KiDS-1000 et DES~Y3.
\textit{Centre}~: horizon sonore $\rdrag$~; la bande orange
indique la contrainte DESI~BAO.
\textit{Droite}~: variation de l'angle acoustique $\Delta\theta_s$.}
\label{fig:scan}
\end{figure}

Trois résultats se dégagent~:

\textbf{1.~Insensibilité de~$z_{\rm eq}$.} Le redshift d'égalité
$z_{\rm eq}=3402.9$ est indépendant de~$d$, confirmant la cohérence
interne du patch~\camb.

\textbf{2.~Forte variation de~$\rdrag$.} Réduire $d$ de~3.0
à~2.65 produit $\rdrag(d=2.65)=125.2$~Mpc, soit~$-14.9\%$,
incompatible avec les mesures BAO DESI~2024 à très haute
signification.

\textbf{3.~Explosion de~$\sig$.} Pour $d=2.65$~:
$\sig=1.312$ ($+62\%$ par rapport à $\LCDM$), exclu par
KiDS-1000 et DES~Y3 à plusieurs dizaines de~$\sigma$.

\textbf{Conclusion~:} toute déviation constante $d\neq 3$ est
doublement exclue. La déviation doit être confinée aux temps tardifs.

% ──────────────────────────────────────────────────────────────
\section{Modèles à $d(z)$ variable~: variable intermédiaire $X(z)$}
\label{sec:models}
% ──────────────────────────────────────────────────────────────

\subsection{Motivation microscopique}

Les simulations BuP sur $N=6$--$16$ qubits mesurent la dimension
spectrale $d_s\approx 2.61\pm 0.03$, stable en taille de système
et en couplage. L'exposant~$\alpha=1.78$ caractérise la loi
microscopique. L'objectif est d'incorporer ce résultat dans un
modèle cosmologique cohérent.

\subsection{La variable $X(z)$}

Nous introduisons la variable intermédiaire~:
\begin{equation}
  X(z) = X_0\,(1+z)^\beta,
  \label{eq:X}
\end{equation}
représentant un paramètre d'ordre effectif contrôlant l'approche
vers la dimension critique. Le profil de dimension devient~:
\begin{equation}
  d(z) = \DC - \frac{\Delta d}{1 + X(z)^\alpha},
  \qquad \alpha = 1.78 \ \text{(fixé microscopique)}.
  \label{eq:dz}
\end{equation}
\textbf{Statut de~$X(z)$.}
$X(z)$ doit être compris comme une variable effective de granulation grossière,
encodant l'approche vers la dimension critique, et non comme un ratio
microscopique directement observable. Sa dynamique $X(z)=X_0(1+z)^\beta$
est phénoménologique~; sa dérivation depuis la dynamique du réseau d'intrication
BuP constitue un objectif de recherche en cours.

Propriétés asymptotiques~:
\begin{align}
  z\to\infty &: \quad X\to\infty \implies d\to \DC \approx 3.06
  \quad \text{(CMB standard ✓)}, \notag\\
  z=0 &: \quad X=X_0 \implies d = \DC - \frac{\Delta d}{1+X_0^\alpha}
  < \DC \quad \text{(déviation tardive ✓)}.
  \notag
\end{align}
Les constantes BuP sont~: $\DC=3.059842935509462$,
$A_c=0.7252150458197096$, $\beta_c=1.4134453781512604$,
déterminant la loi $w(d)=-1+A_c\,(\DC-d)^{\beta_c}$.

\subsection{Contraintes physiques}

Les contraintes suivantes sont appliquées aux deux modèles
(Modèle~C de référence et modèle~Paper~2) pour une comparaison
équitable~:
\begin{itemize}
  \item $d(z=0)\geq 2.40$~: cohérence avec SPARC ($2.644\pm 0.295$);
  \item $\beta\leq 2.0$~: transition cosmologique douce;
  \item $H_0\geq 63.0$~: tension~$H_0$ acceptable.
\end{itemize}

\subsection{Résultats statistiques}

Nous ajustons les modèles aux données DESI~2024 BAO~\cite{desi2024},
Pantheon+~\cite{pantheon} et prior~$\rdrag$ Planck. Le
Tableau~\ref{tab:fits} résume les résultats.

\begin{table}[htbp]
\centering
\caption{Comparaison des modèles cosmologiques sur DESI~2024 + Pantheon+.
$\Delta\chi^2$, $\Delta$AIC et $\Delta$BIC calculés par rapport à $\LCDM$.
Les contraintes physiques $d(z=0)\geq 2.40$, $\beta\leq 2.0$,
$H_0\geq 63$ sont appliquées.}
\label{tab:fits}
\smallskip
\begin{tabular}{lcccccc}
\toprule
Modèle & $n_p$ & $\Delta\chi^2$ & $\Delta$AIC & $\Delta$BIC
& $H_0$ & $d(z{=}0)$ \\
\midrule
$\LCDM$ (réf.) & 3 & --- & --- & --- & 67.6 & --- \\
Modèle~C ($\lambda=1.4$) & 4 & $-9.32$ & $-7.42$ & $-6.16$
& 63.8 & 2.40 \\
Paper~2 ($\alpha=1.78$ fixé) & 5 & $-10.23$ & $-6.23$ & $-3.71$
& 63.6 & 2.40 \\
\bottomrule
\end{tabular}
\end{table}

\begin{figure}[htbp]
\centering
\includegraphics[width=\textwidth]{fig2_models_coherence}
\caption{\textit{Gauche}~: profils $d(z)$ du Modèle~C (Paper~1,
tirets) et Paper~2 (trait plein). \textit{Centre}~: équation
d'état $w(z)$ correspondante. \textit{Droite}~: cohérence
à trois échelles --- simulations BuP ($d_s=2.61\pm 0.03$),
SPARC ($2.644\pm 0.295$, 175~galaxies) et données cosmologiques
($2.40$--$2.71$), sur douze ordres de grandeur.}
\label{fig:coherence}
\end{figure}

Le Tableau~\ref{tab:profile} montre le profil~$d(z)$ best-fit.
En $z=1090$, $d$ est indiscernable de~$\DC$~: les observables
CMB sont inchangées par construction.

\begin{table}[htbp]
\centering
\caption{Profil $d(z)$ pour le modèle Paper~2
($X_0=0.537$, $\beta=2.00$, $\Delta d=0.878$, $\alpha=1.78$).}
\label{tab:profile}
\smallskip
\begin{tabular}{cccc}
\toprule
$z$ & $d(z)$ & $w(z)$ & Époque \\
\midrule
0.0  & 2.400 & $-0.597$ & Aujourd'hui \\
0.3  & 2.583 & $-0.745$ & Pic SNe~Ia \\
1.0  & 2.881 & $-0.936$ & Pic BAO DESI \\
2.0  & 3.010 & $-0.990$ & --- \\
5.0  & 3.055 & $-0.997$ & --- \\
1090 & 3.060 & $-1.000$ & Recombinaison \\
\bottomrule
\end{tabular}
\end{table}

\textbf{Note sur l'AIC.} Paper~2 compte formellement 5~paramètres
contre~4 pour le Modèle~C. Cependant, $\alpha=1.78$ est contraint de
façon indépendante par les simulations BuP --- ce n'est pas un paramètre
libre cosmologique. La pénalité AIC est donc conservative~: l'avantage
statistique de Paper~2 ($\Delta\mathrm{AIC}=-6.23$) est en réalité
sous-estimé, et l'avantage structurel (paramètre microscopique fixé)
est la contribution scientifique centrale.

\subsection{Tension $H_0$}

Les modèles BuP préfèrent $H_0\approx 63$--$64$~km/s/Mpc, en
tension avec Planck ($67.4\pm 0.5$). Cette tension est structurelle
et partagée avec tous les modèles d'énergie noire dynamique fittés
sur BAO+SNe. L'analyse~\camb\ montre que modifier la Friedmann
avec $d(z)$ déplace $H_0^{\rm CMB}$ dans la mauvaise direction.
La résolution nécessite un traitement cohérent des perturbations,
abordé dans la section suivante.

% ──────────────────────────────────────────────────────────────
\section{$G_{\rm eff}$ cohérent dans les perturbations~:
  résolution de la tension~$\sigma_8$}
\label{sec:sigma8}
% ──────────────────────────────────────────────────────────────

\subsection{Le problème du traitement incomplet}

Dans Paper~2, $\Geff=\frac{2}{d(z)-1}\,G$ est appliqué à
l'équation de Friedmann (fond) uniquement. Les perturbations
utilisent~$G$ standard. Cette inconsistance --- modifier
l'expansion sans modifier la croissance --- produit $\sig\approx 0.86$
(tension KiDS $\approx 4.7\sigma$). C'est un artefact du traitement
incomplet, pas une prédiction physique.

\subsection{Traitement cohérent}

L'équation de croissance linéaire avec $\Geff(z)$ cohérent depuis
la recombinaison~:
\begin{equation}
  \ddot{D} + \left(\frac{3}{a} + \frac{d\ln E}{da}\right)\dot{D}
  = \frac{3\,\Omega_m}{2\,a^5\,E^2}\,\frac{\Geff(z)}{G}\,D,
  \label{eq:growth}
\end{equation}
où $E(z)=H(z)/H_0$ et $\Geff(z)/G=2/(d(z)-1)$. La normalisation
est $\sig=0.8114\times D_{\rm BuP}/D_{\LCDM}$.

\subsection{Mécanisme géométrique}

Le Tableau~\ref{tab:geff} et la Figure~\ref{fig:sigma8} illustrent
la variation de~$\Geff$ avec le redshift.

\begin{table}[htbp]
\centering
\caption{Profil $\Geff(z)/G$ pour le best-fit Paper~2.
La clé du mécanisme est la valeur $\DC=3.0598>3$~: à haut
redshift, $\Geff\approx 0.97\,G$ (légère inhibition).}
\label{tab:geff}
\smallskip
\begin{tabular}{cccc}
\toprule
$z$ & $d(z)$ & $\Geff/G$ & Effet \\
\midrule
0.0   & 2.400 & 1.429 & $+42.9\%$ (renforcement) \\
0.5   & 2.694 & 1.181 & $+18.1\%$ (renforcement) \\
1.0   & 2.881 & 1.063 & $+6.3\%$  (renforcement) \\
2.0   & 3.010 & 0.995 & $-0.5\%$  (inhibition) \\
5.0   & 3.055 & 0.973 & $-2.7\%$  (inhibition) \\
1090  & 3.060 & 0.971 & $-2.9\%$  (inhibition) \\
\bottomrule
\end{tabular}
\end{table}

Le mécanisme est le suivant~: $\sig$ intègre la croissance
depuis $z=1000$. La légère inhibition à haut~$z$ ($\sim 3\%$
sur~$\sim 99\%$ de l'histoire cosmique) domine le renforcement
à bas~$z$ ($\sim 43\%$ sur~$\sim 1\%$ de l'histoire). Ce
mécanisme est entièrement géométrique --- il découle de
$\DC>3$, propriété fondamentale du réseau d'intrication BuP,
sans paramètre libre supplémentaire.

\subsection{Résultats}

\begin{table}[htbp]
\centering
\caption{$\sigma_8$ en fonction de~$\Delta d$ (traitement cohérent
$\Geff$, $X_0=0.537$, $\beta=2.0$, $\Omega_m=0.3448$ fixés).
Toutes les valeurs sont compatibles avec KiDS-1000 à moins de
$1\sigma$.}
\label{tab:sigma8}
\smallskip
\begin{tabular}{cccc}
\toprule
$\Delta d$ & $d(z{=}0)$ & $\sig$ & Tension KiDS \\
\midrule
0.300 & 2.834 & 0.751 & $-0.8\sigma$ \\
0.500 & 2.684 & 0.757 & $-0.4\sigma$ \\
0.650 & 2.555 & 0.763 & $-0.2\sigma$ \\
\textbf{0.878} & \textbf{2.400} & \textbf{0.772} & $+0.3\sigma$ \\
0.900 & 2.384 & 0.773 & $+0.3\sigma$ \\
\bottomrule
\end{tabular}
\end{table}

\begin{figure}[htbp]
\centering
\includegraphics[width=\textwidth]{fig3_sigma8_geff}
\caption{\textit{Gauche}~: profil $\Geff(z)/G$ pour le best-fit
Paper~2. La zone bleue indique l'inhibition à haut~$z$~;
la zone rouge le renforcement à bas~$z$.
\textit{Droite}~: $\sig$ en fonction de $d(z=0)$, avec
traitement cohérent de~$\Geff$. Les bandes colorées indiquent
les contraintes KiDS-1000 et DES~Y3.}
\label{fig:sigma8}
\end{figure}

Pour le best-fit Paper~2~:
\begin{equation}
  \sig^{\rm BuP} = 0.772, \qquad
  \text{tension KiDS-1000}~: +0.3\sigma.
\end{equation}
Ce résultat est robuste sur toute la gamme physique
$\Delta d\in[0.3,\,0.9]$.

\subsection{Tableau de cohérence}

\begin{table}[htbp]
\centering
\caption{Cohérence du cadre BuP avec les données actuelles.}
\label{tab:coherence}
\smallskip
\begin{tabular}{lllll}
\toprule
Observable & Valeur BuP & Données & Tension & Statut \\
\midrule
$\rdrag$ (BAO) & inchangé & DESI~2024 & $<1\sigma$ & ✓ \\
$\sig$ & 0.772 & KiDS $0.766\pm 0.020$ & $+0.3\sigma$ & ✓ \\
$d(z{=}0)$ & 2.40 & SPARC $2.644\pm 0.295$ & $+0.8\sigma$ & ✓ \\
$\Delta$BIC & $-3.71$ & DESI+Pantheon+ & signal & ✓ \\
CMB ($\theta_s$) & inchangé & Planck~2018 & $<0.1\sigma$ & ✓ \\
$H_0$ & 63.7 & Planck $67.4\pm 0.5$ & $3.7\sigma$ & ⚠ \\
\bottomrule
\end{tabular}
\end{table}

\subsection{Limites}

Le calcul de~$\sig$ est une intégration de l'équation de croissance
linéaire (ODE), pas un calcul Boltzmann complet. Il n'inclut pas le
transfert de puissance $k$-dépendant, les effets non-linéaires, ni la
rétroaction des neutrinos massifs avec pleine précision. Ce résultat
devrait être confirmé par une implémentation complète dans un code Boltzmann
(\camb\ ou \textsc{class})~; une telle implémentation via une classe
d'énergie sombre personnalisée est en cours de développement et constituera
l'objet d'une publication séparée.

% ──────────────────────────────────────────────────────────────
\section{Discussion}
\label{sec:discussion}
% ──────────────────────────────────────────────────────────────

\subsection{Cohérence multi-échelle}

La valeur $d_s\approx 2.61\pm 0.03$ est cohérente entre~:
(i)~les simulations BuP sur $N=6$--$16$ qubits (niveau microscopique,
échelle $\sim 10^{-2}$~kpc)~;
(ii)~les courbes de rotation galactiques SPARC, 175~galaxies
($d_{\rm eff}=2.644\pm 0.295$, échelle $\sim 10^4$~kpc)~;
(iii)~les données cosmologiques DESI+Pantheon+ ($d(z=0)\approx 2.40$--$2.71$,
échelle $\sim 10^{10}$~kpc).
Cette cohérence sur douze ordres de grandeur --- non imposée ---
constitue l'argument observationnel le plus fort en faveur du cadre BuP.

\subsection{Rôle de $\DC>3$}

La valeur $\DC=3.0598>3$ joue un rôle doublement crucial~:
(i)~elle garantit que $d\to\DC$ à haut~$z$, préservant
les observables du CMB~;
(ii)~elle implique $\Geff\approx 0.97\,G$ à haut~$z$, inhibant
légèrement la croissance des structures et résolvant naturellement
la tension~$\sig$. Cette valeur est déterminée par la dynamique
microscopique du réseau d'intrication, non par un ajustement cosmologique.

\subsection{Tension $H_0$}

La tension $H_0\approx 3.7\sigma$ (BAO+SNe vs CMB) est structurelle
dans le cadre actuel~: $H_0\approx 63.7$~km/s/Mpc dans BuP résulte
directement de $d(z=0)\approx 2.40$, qui ralentit l'expansion à bas~$z$
et requiert un~$H_0$ plus faible pour reproduire les distances aux supernovas.
Cette tension n'est pas spécifique à BuP~; elle est partagée avec tous les
modèles d'énergie sombre dynamique contraints sur ces mêmes données~\cite{tensions2025}.
La résolution simultanée des tensions $H_0$ et~$\sig$ dans un cadre cohérent
nécessite probablement soit une modification de la physique de recombinaison
(~$r_{\rm drag}$ plus grand), soit un couplage matière-géométrie dépendant
de l'échelle~$k$.

\subsection{Perspectives}

Les prochaines étapes incluent~: (a)~l'implémentation complète
de~$\Geff(z,k)$ dans \camb\ via une classe d'énergie sombre
personnalisée~; (b)~la contrainte simultanée sur le spectre CMB
avec~$d$ libre~; (c)~l'analyse de taille finie FSS sur $N=20$--$32$
qubits pour faire converger~$d_s$ vers la plage cosmologique~;
(d)~le test sur JWST~: l'excès de galaxies massives à $z>5$ prédit
par la formation accélérée des structures dans BuP à bas~$z$.

% ──────────────────────────────────────────────────────────────
\section{Conclusion}
\label{sec:conclusion}
% ──────────────────────────────────────────────────────────────

Nous avons présenté un cadre cohérent pour la cosmologie BuP,
articulé en trois résultats principaux~:

\textbf{(i)~Exclusion robuste de $d\neq 3$ constant.} Le scan \camb\
montre que $\rdrag$ et~$\sig$ sont tous deux incompatibles avec
$d\neq 3$ constant, de façon indépendante.

\textbf{(ii)~Signal cosmologique tardif.} Le modèle Paper~2,
$d(z)=\DC-\Delta d/(1+X(z)^\alpha)$ avec $\alpha=1.78$ fixé
par les simulations microscopiques, donne $\Delta\mathrm{AIC}=-6.23$
et $\Delta\mathrm{BIC}=-3.71$ sur DESI~2024+Pantheon+. La cohérence
entre simulations BuP, SPARC et données cosmologiques sur douze
ordres de grandeur constitue le résultat le plus fort.

\textbf{(iii)~Résolution de la tension~$\sig$.} L'implémentation
cohérente de $\Geff(z)=\frac{2}{d(z)-1}\,G$ dans les équations
de perturbation donne $\sig=0.772$, en accord avec KiDS-1000
à~$0.3\sigma$. Le mécanisme est géométrique~: la légère
sur-dimensionnalité de l'Univers primordial ($\DC>3$) inhibe
la croissance à haut redshift et compense le renforcement tardif.

Une réalisation complète du programme requiert~: (a)~l'implémentation
Boltzmann de~$\Geff(z)$~; (b)~la dérivation de~$X(z)$ depuis la
dynamique du réseau d'intrication BuP~; (c)~l'ajustement simultané
aux spectres CMB avec~$d$ libre.

La légère sur-dimensionnalité de l'Univers primordial ($\DC>3$) fournit
une explication géométrique unifiée pour l'accélération cosmique tardive
et la suppression de la croissance des structures --- sans paramètre libre
au-delà de ceux contraints par les simulations microscopiques.
Ce résultat, robuste sur douze ordres de grandeur, suggère que la dimension
effective de l'espace-temps est une propriété émergente mesurable, et non
un postulat fixé~à~priori.

% ──────────────────────────────────────────────────────────────
\bibliographystyle{plain}
\begin{thebibliography}{99}

\bibitem{hamdad2026bup}
F.~Hamdad,
\textit{Bottom-Up Quantum Gravity~: emergent space, time and gravity from
quantum entanglement}, 2026,
\url{https://github.com/Farid-Hamdad/Bottom-up-Quantum-Gravity}.

\bibitem{desi2024}
DESI Collaboration,
\textit{DESI~2024 VI~: Cosmological Constraints from Baryon Acoustic Oscillations},
arXiv:2404.03002 (2024).

\bibitem{pantheon}
D.~Brout \textit{et al.},
\textit{The Pantheon+ Analysis~: Cosmological Constraints},
ApJ \textbf{938}, 110 (2022).

\bibitem{kids}
KiDS Collaboration,
\textit{KiDS-1000 Cosmology~: Cosmic shear constraints},
A\&A \textbf{645}, A104 (2021).

\bibitem{des}
DES Collaboration,
\textit{Dark Energy Survey Year~3 results~: Cosmological constraints},
Phys.\ Rev.\ D \textbf{105}, 023520 (2022).

\bibitem{planck2018}
Planck Collaboration,
\textit{Planck 2018 results VI. Cosmological parameters},
A\&A \textbf{641}, A6 (2020).

\bibitem{tensions2025}
E.~Di~Valentino \textit{et al.},
\textit{Tensions in Cosmology 2025},
Nature Astronomy (2026), \url{https://doi.org/10.1038/s41550-026-02781-1}.

\end{thebibliography}

\end{document}

